% !PRINTED 70
\weberpage{70}
The following theorem now holds:

\gesperrt{3. If $\Omega'$ and $\Omega''$ are two divisors of $\Omega_m$ which belong to the
groups $\mathfrak{A}'$ and $\mathfrak{A}''$, then $\Omega''$ is a divisor
of $\Omega'$ if and only if $\mathfrak{A}'$ is a divisor
of $\mathfrak{A}''$.}

For if $\Omega''$ is a divisor of $\Omega'$, then all numbers
of $\Omega''$ must admit the substitutions of the group $\mathfrak{A}'$ to which $\Omega'$ belongs; hence $\mathfrak{A}'$ must be contained in $\mathfrak{A}''$. And if conversely $\mathfrak{A}'$ is contained in $\mathfrak{A}''$,
then all numbers of $\Omega''$ admit
all the substitutions of $\mathfrak{A}'$ and are therefore rationally
representable by a primitive number of the field $\Omega'$; that is, $\Omega''$
is contained in $\Omega'$ (§. 144 of the first volume).

The field $\Omega_m$ contains as divisors all the fields $\Omega_{m_1}$ for
which $m_1$ is a divisor of $m$; for $\Omega_{m_1}$ consists of all
rational functions of $r^{\frac{m}{m_1}}$. Now we have
\[
r^{a\frac{m}{m_1}} = r^{\frac{m}{m_1}},
\]
if and only if
\weq{4}{a \equiv 1\,(\mathrm{mod}\ m_1).}

The numbers $a$ in $\mathfrak{N}$ that satisfy the congruence (4) thus
form the group to which the field $\Omega_{m_1}$ belongs. We will
denote this group by $\mathfrak{A}_{m_1}$ and symbolically put
\weq{5}{\mathfrak{A}_{m_1} \equiv 1\,(\mathrm{mod}\ m_1).}

Since the degree of the field $\Omega_{m_1}$ is equal to $\varphi(m_1)$, which is
at the same time the index of the divisor $\mathfrak{A}_{m_1}$ of $\mathfrak{N}$, the degree
of $\mathfrak{A}_{m_1}$ is equal to $\varphi(m):\varphi(m_1)$.

If $m = m_1 m_2$ and $m_1$ is relatively prime to $m_2$, one obtains
all the numbers of $\mathfrak{A}_{m_1}$ from $a = 1 + x m_1$, if $x$ is determined
so that $a$ becomes congruent modulo $m_2$ with all the numbers relatively prime to
$m_2$. From this there results, in agreement
with the above general determination, the degree of $\mathfrak{A}_{m_1}$ for
this case $=\varphi(m_2)$.

But if now $m_1$ and $m_2$ are any two divisors of $m$, and
$\mu$ is the greatest common divisor of $m_1$ and $m_2$, then
the greatest common divisor of the two fields $\Omega_{m_1}, \Omega_{m_2}$
will belong to the group which one may regard as the least common
multiple of the two groups
\weq{6}{\mathfrak{A}_{m_1} \equiv 1\,(\mathrm{mod}\ m_1),\quad \mathfrak{A}_{m_2} \equiv 1\,(\mathrm{mod}\ m_2)}
%CONT>>
